\documentclass[../../main.tex]{subfiles}% 注意这里的文件路径不能用 ./main.tex ,否则用latexmk编译子文件会报错
\graphicspath{{\subfix{./image/}}{\subfix{../../image/}}} % 指定图片引用路径,后续可以直接使用图片文件名
% 如果图片存放在上级目录,则使用 ../../image/ 作为备用路径

% 例如:
% \begin{figure}[H]
% \centering
% \includegraphics[width=0.3\textwidth]{图.png}
% \caption{}
% \label{figure:图}
% \end{figure}
% 注意:上述\label{}一定要放在\caption{}之后,否则引用图片序号会只会显示??.

\begin{document}

\section{习题}

\begin{example}
用特征线法求解下列 Cauchy 初值问题：
\begin{enumerate}[(1)]
\item $\begin{cases}
u_t+2u_x=0, & (x,t)\in\mathbb{R}_+^2, \\
u(x,0)=x^2, & x\in\mathbb{R};
\end{cases}$
\item $\begin{cases}
u_t+2u_x+u=xt, & (x,t)\in\mathbb{R}_+^2, \\
u(x,0)=2-x, & x\in\mathbb{R};
\end{cases}$
\item $\begin{cases}
2u_t-u_x+xu=0, & (x,t)\in\mathbb{R}_+^2, \\
u(x,0)=2xe^{x^2/2}, & x\in\mathbb{R};
\end{cases}$
\item $\begin{cases}
u_t+(1+x^2)u_x-u=0, & (x,t)\in\mathbb{R}_+^2, \\
u(x,0)=\arctan x, & x\in\mathbb{R};
\end{cases}$
\item $\begin{cases}
u_t+A\cdot Du+cu=0, & (x,t)\in\mathbb{R}_+^{n+1}, \\
u(x,0)=\varphi(x), & x\in\mathbb{R}^n,
\end{cases}$
这里 $A\in\mathbb{R}^n$ 是常向量，$c\in\mathbb{R}$ 是常数。
\end{enumerate}
\end{example}
\begin{solution}
\begin{enumerate}[(1)]
\item 方程的特征线方程为$\frac{\mathrm{d}x}{\mathrm{d}t}=2$, 解得过点$(0,\xi)$的特征线为
\begin{align*}
x(t)=\xi+2t,\quad t\in\mathbb{R}.
\end{align*}
沿特征线, 原方程化为
\begin{align*}
\frac{\mathrm{d}}{\mathrm{d}t}u(x(t),t)=u_t+2u_x=0\,(\forall t\in\mathbb{R}) \Longrightarrow u(x(t),t)=f(\xi)=f(x(t)-2t)\,(\forall t\in\mathbb{R}).
\end{align*}
其中$f$为一元可微函数. 由$\xi$的任意性得
\begin{align*}
u(x,t)=f(x-2t),\quad \forall (x,t)\in\mathbb{R}^2.
\end{align*}
将$t=0,u(x,0)=x^2$代入上式得$f(x)=x^2$. 故
\begin{align*}
u(x,t)=(x-2t)^2.
\end{align*}

\item 
方程的特征线方程为$\frac{\mathrm{d}x}{\mathrm{d}t}=2$, 解得过点$(0,\xi)$的特征线为
\begin{align*}
x(t)=\xi+2t,\quad t\in\mathbb{R}.
\end{align*}
沿特征线, 原方程化为
\begin{align*}
\frac{\mathrm{d}}{\mathrm{d}t}u(x(t),t)=u_t+2u_x=x(t)\cdot t-u(x(t),t)=t(\xi+2t)-u(x(t),t).
\end{align*}
由一阶线性常微分方程的常数变易法, 解得
\begin{align*}
u(x(t),t)=\xi(t-1)+2(t^2-2t+2)+Ce^{-t}=(x(t)-2t)(t-1)+2(t^2-2t+2)+f(\xi)e^{-t},
\end{align*}
其中$f$为一元可微函数. 由$\xi$的任意性可得
\begin{align*}
u(x,t)=(x-2t)(t-1)+2(t^2-2t+2)+Ce^{-t},\quad \forall (x,t)\in\mathbb{R}^2.
\end{align*}
将$t=0,u(x,0)=2-x$代入上式得$f\equiv -2$. 故
\begin{align*}
u(x,t)=(x-2t)(t-1)+2(t^2-2t+2)-2e^{-t}=x(t-1)-2t+4-2e^{-t}.
\end{align*}

\item 化原方程为
\begin{align*}
u_t-\frac12 u_x+\frac{x}{2}u=0.
\end{align*}
特征线方程为
\begin{align*}
\frac{\mathrm{d}x}{\mathrm{d}t}=-\frac12,
\end{align*}
解得过点$(0,\xi)$的特征线为
\begin{align*}
x(t)=\xi-\frac{t}{2},\quad t\in\mathbb{R}.
\end{align*}
沿此特征线，有
\begin{align*}
\frac{\mathrm{d}}{\mathrm{d}t}u(x(t),t)=u_t+u_x\frac{\mathrm{d}x}{\mathrm{d}t}
=u_t-\frac12 u_x=-\frac{x}{2}u.
\end{align*}
代入$x(t)=\xi-\frac{t}{2}$，得
\begin{align*}
\frac{\mathrm{d}u}{\mathrm{d}t}=-\frac{1}{2}\left(\xi-\frac{t}{2}\right)u
=\left(-\frac{\xi}{2}+\frac{t}{4}\right)u.
\end{align*}
分离变量并积分：
\begin{align*}
&\frac{\mathrm{d}u}{u}=\left(-\frac{\xi}{2}+\frac{t}{4}\right)dt
\Longrightarrow
\ln|u|=-\frac{\xi t}{2}+\frac{t^2}{8}+f(\xi)\\
\Longrightarrow &u(x(t),t)=f(\xi)\exp\left(-\frac{\xi t}{2}+\frac{t^2}{8}\right)=f(x(t)+\frac{t}{2})\exp\left(-\frac{(x(t)+\frac{t}{2}) t}{2}+\frac{t^2}{8}\right),
\end{align*}
其中$f$为一元可微函数. 由$\xi$的任意性可得
\begin{align*}
u(x,t)=f\!\left(x+\frac{t}{2}\right)\exp\left(-\frac{(x+\frac{t}{2})t}{2}+\frac{t^2}{8}\right),\quad \forall (x,t)\in \mathbb{R}^2.
\end{align*}
代入初值$t=0,u(x,0)=2xe^{x^2/2}$得
$f(x)=2\left(x+\frac{t}{2}\right)e^{\frac{\left(x+\frac{t}{2}\right)^2}{2}}.$
因此
\begin{align*}
u(x,t)=(2x+t)e^{x^2/2}.
\end{align*}

\item 原方程可写为
\begin{align*}
u_t+(1+x^2)u_x=u.
\end{align*}
其特征线方程为
\begin{align*}
\frac{\mathrm{d}x}{\mathrm{d}t}=1+x^2,
\end{align*}
解得过点$(0,\xi)$的特征线满足
\begin{align*}
\arctan x - t = \arctan \xi,
\end{align*}
即
\begin{align*}
\xi=\tan(\arctan x - t).
\end{align*}
沿特征线，有
\begin{align*}
\frac{\mathrm{d}}{\mathrm{d}t}u(x(t),t)=u_t+u_x\frac{\mathrm{d}x}{\mathrm{d}t}
=u_t+(1+x^2)u_x=u.
\end{align*}
因此
\begin{align*}
\frac{\mathrm{d}u}{\mathrm{d}t}=u \;\Longrightarrow\; u(x(t),t)=f(\xi) e^t=f(\tan(\arctan x(t) - t))e^t,
\end{align*}
其中$f$为一元可微函数. 由$\xi$的任意性可得
\begin{align*}
u(x,t)=f(\tan(\arctan x - t))e^t,\quad \forall (x,t)\in \mathbb{R}^2.
\end{align*}
代入初值$t=0,u(x,0)=\arctan x$得$f(x)=\arctan x.$故
\begin{align*}
u(x,t)=(\arctan x - t)e^t.
\end{align*}

\item 特征线满足向量常微分方程
\begin{align*}
\frac{\mathrm{d}x}{\mathrm{d}t}=A,
\end{align*}
其过点$(0,\xi)$的解为
\begin{align*}
X(t)=\xi+At,\quad t\in\mathbb{R}^n,
\end{align*}
沿此特征线，有
\begin{align*}
\frac{\mathrm{d}}{\mathrm{d}t}u(X(t),t)=u_t+\nabla_x u\cdot \frac{\mathrm{d}x}{\mathrm{d}t}
=u_t+A\cdot Du=-cu(X(t),t)\Longrightarrow u(X(t),t)=f(\xi)e^{-ct}=f(X(t)-At)e^{-ct}.
\end{align*}
由$\xi $的任意性可得
\begin{align*}
u(x,t)=f(x-At)e^{-ct},\quad \forall (x,t)\in \mathbb{R}^n\times \mathbb{R}^n.
\end{align*}
代入初值$u(0)=u(x,0)=\varphi(x)$得
$f(x)=\varphi(x)$.进而
\begin{align*}
u(x,t)=e^{-ct}\varphi(x-At).
\end{align*}
\end{enumerate}
\end{solution}

\begin{example}
用特征线法求解拟线性方程的 Cauchy 初值问题
\begin{align*}
\begin{cases}
u_t+uu_x=0, & (x,t)\in\mathbb{R}_+^2, \\
u(x,0)=\varphi(x), & x\in\mathbb{R},
\end{cases}
\end{align*}
其中 $\varphi(x)$ 是一个递增函数。（提示：考虑特征线 $\frac{\mathrm{d}x}{\mathrm{d}t}=u$）
\end{example}

\begin{proposition}
\begin{enumerate}[(1)]
\item 设 $u=u(x,y)$ 在凸连通区域 $\Omega\subset\mathbb{R}^2$ 上满足方程 $u_{xy}=0$，求此方程的所有解。
\item 设 $a$ 为正常数。利用变量替换 $\xi=x+at,\ \eta=x-at$ 证明 $u_{tt}-a^2u_{xx}=0$ 当且仅当 $u_{\xi\eta}=0$。
\item 利用 (1), (2) 推导 D'Alembert 公式。
\end{enumerate}
\end{proposition}
\begin{proof}
\begin{enumerate}[(1)]
\item 

\item 

\item 
\end{enumerate}
\end{proof}

\begin{example}
设 $ABCD$ 表示 $\mathbb{R}_+^2$ 上的一个平行四边形，其中两条边 $AB,CD$ 平行于波动方程 $u_{tt}-a^2u_{xx}=0$ ($a>0$) 的一条特征线 $x-at=0$，另两条边 $AD,BC$ 平行于另一条特征线 $x+at=0$。这样的平行四边形称为特征平行四边形。设 $u(x,t)$ 是波动方程的解。证明：
\begin{align*}
u(A)+u(C)=u(B)+u(D),
\end{align*}
其中 $u(A),u(B),u(C),u(D)$ 分别表示函数 $u$ 在顶点 $A,B,C,D$ 处的函数值。
\end{example}

\begin{example}
假设电场强度 $\bm{E}=(E_1,E_2,E_3)$ 和磁场强度 $\bm{B}=(B_1,B_2,B_3)$ 满足 Maxwell 方程组
\begin{align*}
\begin{cases}
\frac{1}{c}\cdot\frac{\partial \bm{E}}{\partial t}=\operatorname{curl}\bm{B}, \\
\frac{1}{c}\cdot\frac{\partial \bm{B}}{\partial t}=-\operatorname{curl}\bm{E}, \\
\operatorname{div}\bm{E}=\operatorname{div}\bm{B}=0,
\end{cases}
\end{align*}
这里 $c$ 为光速。证明：它们的分量 $E_1,E_2,E_3,B_1,B_2,B_3$ 满足波动方程 $u_{tt}-c^2\Delta u=0$。
\end{example}

\begin{example}
证明：方程
\begin{align*}
\left(1-\frac{x}{h}\right)^2\frac{\partial^2u}{\partial t^2}
=a^2\frac{\partial}{\partial x}\left[\left(1-\frac{x}{h}\right)^2\frac{\partial u}{\partial x}\right]
\end{align*}
的通解可以表示成
\begin{align*}
u(x,t)=\frac{F(x-at)+G(x+at)}{h-x},
\end{align*}
其中 $F,G$ 是任意二次连续可微函数，$h>0,\ a>0$ 为常数。（提示：考虑函数替换 $v=(h-x)u$）
\end{example}

\begin{example}
直接用特征线法来求解初值问题 \eqref{eq:wave-cauchy1d}。
\end{example}

\begin{example}
直接验证定理 \eqref{eq:wave-cauchy1d} 和推论。
\end{example}

\begin{example}
当初值 $u(x,0)=\varphi(x),\ u_t(x,0)=\psi(x)$ 满足什么条件时，一维齐次波动方程的初值问题的解仅由右行波组成（即通解为 $G(x-at)$ 的形式）？
\end{example}

\begin{example}
已知初值问题
\begin{align*}
\begin{cases}
u_{tt}-u_{xx}=0, & x\in\mathbb{R},\ t>a, \\
u|_{t=a}=u_0(x), & x\in\mathbb{R}, \\
u_t|_{t=a}=u_1(x), & x\in\mathbb{R},
\end{cases}
\end{align*}
其中 $a\neq\pm 1$。若初值在 $b\leqslant x\leqslant c$ 上给定，试问它的解在什么区域确定。
\end{example}

\begin{example}
证明：初值问题
\begin{align*}
\begin{cases}
u_{tt}-u_{xx}=6(x+t), & x\in\mathbb{R},\ t>x, \\
u|_{t=x}=0, & x\in\mathbb{R}, \\
u_t|_{t=x}=\psi(x), & x\in\mathbb{R}
\end{cases}
\end{align*}
有解的充分必要条件是 $\psi(x)-3x^2=\text{const}$；若上述问题有解，则问题的解不是唯一的。试问：当在直线 $t=ax$ 上给定初值时，为什么在 $a=\pm 1$ 和 $a\neq\pm 1$ 的情形，关于解的存在性唯一性的结论完全不同？
\end{example}

\begin{example}
设 $a>0$ 为常数。利用 Fourier 变换求解三维波动方程初值问题
\begin{align*}
\begin{cases}
u_{tt}-a^2(u_{xx}+u_{yy}+u_{zz})=f(x,y,z,t), & (x,y,z,t)\in\mathbb{R}^4, \\
u|_{t=0}=\varphi(x,y,z), & (x,y,z)\in\mathbb{R}^3, \\
u_t|_{t=0}=\psi(x,y,z), & (x,y,z)\in\mathbb{R}^3.
\end{cases}
\end{align*}
\end{example}

\begin{example}
设 $a>0$ 为常数。利用 Fourier 变换求解二维波动方程初值问题
\begin{align*}
\begin{cases}
u_{tt}-a^2(u_{xx}+u_{yy})=f(x,y,t), & (x,y,t)\in\mathbb{R}_+^3, \\
u|_{t=0}=\varphi(x,y), & (x,y)\in\mathbb{R}^2, \\
u_t|_{t=0}=\psi(x,y), & (x,y)\in\mathbb{R}^2.
\end{cases}
\end{align*}
\end{example}

\begin{example}
设 $u=u(x,y,z,t)$ 是三维波动方程初值问题
\begin{align*}
\begin{cases}
u_{tt}-a^2(u_{xx}+u_{yy}+u_{zz})=0, & (x,y,z,t)\in\mathbb{R}_+^4, \\
u|_{t=0}=f(x)+g(y), & (x,y,z)\in\mathbb{R}^3, \\
u_t|_{t=0}=\varphi(y)+\psi(z), & (x,y,z)\in\mathbb{R}^3
\end{cases}
\end{align*}
的解，其中 $a>0$ 为常数。求解 $u$ 的表达式。
\end{example}

\begin{example}
设函数 $\Phi\in C^2(\mathbb{R})$，向量 $\alpha\in\mathbb{R}^n$ 且 $|\alpha|=1$，则 $\Phi(\alpha\cdot x+at)$ 满足 $n$ 维波动方程
\begin{align*}
u_{tt}-a^2\Delta u=0,
\end{align*}
这里 $a>0$ 为常数，$x=(x_1,\cdots,x_n)$，$\Delta=\frac{\partial^2}{\partial x_1^2}+\cdots+\frac{\partial^2}{\partial x_n^2}$。波动方程的这种形式的解称为平面波解。
\end{example}

\begin{example}
设 $a$ 为正常数。证明：三维波动方程
\begin{align*}
u_{tt}-a^2\Delta u=0,\quad (x,t)\in\mathbb{R}_+^4
\end{align*}
的所有径向对称的解可以表示为
\begin{align*}
u(x,t)=\frac{F(|x|-at)+G(|x|+at)}{|x|},\quad x\neq 0.
\end{align*}
\end{example}

\begin{example}
设 $a$ 为正常数。求解下列多维初值问题：
\begin{enumerate}[(1)]
\item $\begin{cases}
u_{tt}-a^2(u_{xx}+u_{yy})=0, & (x,y,t)\in\mathbb{R}_+^3, \\
u|_{t=0}=x^2(x+y),\quad u_t|_{t=0}=0, & (x,y)\in\mathbb{R}^2;
\end{cases}$
\item $\begin{cases}
u_{tt}-a^2(u_{xx}+u_{yy}+u_{zz})=0, & (x,y,z,t)\in\mathbb{R}_+^4, \\
u|_{t=0}=x^2+y^2z,\quad u_t|_{t=0}=1+y, & (x,y,z)\in\mathbb{R}^3.
\end{cases}$
\end{enumerate}
\end{example}

\begin{example}
设 $a$ 为正常数，函数 $\varphi(x),\psi(x)\in C^2[0,+\infty)$，$g(t)\in C^2[0,+\infty)$ 满足相容性条件
\begin{align*}
g(0)=\varphi(0),\quad g'(0)=\psi(0),\quad g''(0)=a^2\varphi''(0).
\end{align*}
试求解半无界问题
\begin{align*}
\begin{cases}
u_{tt}-a^2u_{xx}=0, & (x,t)\in\mathbb{R}_+\times\mathbb{R}_+, \\
u(x,0)=\varphi(x),\quad u_t(x,0)=\psi(x), & x\geqslant 0, \\
u(0,t)=g(t), & t\geqslant 0.
\end{cases}
\end{align*}
\end{example}

\begin{example}
试给出半无界问题
\begin{align*}
\begin{cases}
u_{tt}-a^2u_{xx}=f(x,t), & (x,t)\in\mathbb{R}_+\times\mathbb{R}_+, \\
u(x,0)=\varphi(x),\quad u_t(x,0)=\psi(x), & x\geqslant 0, \\
u_x(0,t)=g(t), & t\geqslant 0
\end{cases}
\end{align*}
有古典解的相容性条件。
\end{example}

\begin{example}
利用惟一性结果直接证明：当初值 $\varphi(x),\psi(x)$ 和非齐次项 $f(x,t)$ 是关于 $x$ 的偶函数时初值问题 \eqref{eq:wave-cauchy1d} 的解 $u(x,t)$ 也是关于 $x$ 的偶函数。根据以上事实，用延拓法求解半无界问题
\begin{align*}
\begin{cases}
u_{tt}-a^2u_{xx}=f(x,t), & (x,t)\in\mathbb{R}_+\times\mathbb{R}_+, \\
u(x,0)=u_t(x,0)=0, & x\geqslant 0, \\
u_x(0,t)=0, & t\geqslant 0,
\end{cases}
\end{align*}
并说明当 $f(x,t)$ 满足什么条件时，导出的形式解确实是问题的解。这里 $a$ 为正常数。
\end{example}

\begin{example}
利用延拓法求解半无界问题
\begin{align*}
\begin{cases}
u_{tt}-a^2u_{xx}=0, & (x,t)\in\mathbb{R}_+\times\mathbb{R}_+, \\
u(x,0)=u_t(x,0)=0, & x\geqslant 0, \\
u_x(0,t)=A\sin\omega t, & t\geqslant 0,
\end{cases}
\end{align*}
并给出物理解释。这里 $a$ 为正常数，$A$ 为常数。
\end{example}

\begin{example}
设 $u=u(x,y,t)$ 是初值问题
\begin{align*}
\begin{cases}
u_{tt}-4(u_{xx}+u_{yy})=0, & (x,y,t)\in\mathbb{R}_+^3, \\
u(x,y,0)=0, & (x,y)\in\mathbb{R}^2, \\
u_t(x,y,0)=\psi(x,y), & (x,y)\in\mathbb{R}^2
\end{cases}
\end{align*}
的解，其中当 $(x,y)\in\Omega$ 时，$\psi(x,y)\equiv 0$；当 $(x,y)\in\mathbb{R}^2\setminus\Omega$ 时，$\psi(x,y)>0$。这里 $\Omega$ 是正方形 $\{(x,y)\in\mathbb{R}^2: |x|\leqslant 1,\ |y|\leqslant 1\}$。试指出当 $t>0$ 时，$u(x,y,t)\equiv 0$ 的区域。
\end{example}

\begin{example}
设 $u=u(x,t)$ 是半无界问题
\begin{align*}
\begin{cases}
u_{tt}-a^2u_{xx}=0, & (x,t)\in\mathbb{R}_+\times\mathbb{R}_+, \\
u(x,0)=\varphi(x), & x\geqslant 0, \\
u_t(x,0)=\psi(x), & x\geqslant 0, \\
u(0,t)=g(t), & t\geqslant 0
\end{cases}
\end{align*}
的解，其中 $a$ 为正常数，且
\begin{align*}
\varphi(x),\psi(x)=\begin{cases}
0, & x\in[0,1], \\
>0, & x\in(1,+\infty),
\end{cases}\qquad
g(t)=\begin{cases}
0, & t\in[0,1], \\
>0, & t\in(1,+\infty).
\end{cases}
\end{align*}
试指出当 $t>0$ 时，$u(x,t)\equiv 0$ 的区域。
\end{example}

\begin{example}
试问：半无界问题
\begin{align*}
\begin{cases}
u_{tt}-u_{xx}+u_t+u_x=0, & (x,t)\in\mathbb{R}_+\times\mathbb{R}_+, \\
u(x,0)=\varphi(x), & x\geqslant 0, \\
u_t(x,0)=\psi(x), & x\geqslant 0, \\
u(0,t)=0, & t\geqslant 0
\end{cases}
\end{align*}
能否直接用对称开拓法求解？为什么？试用特征线法求解此半无界问题。
提示：利用等式
\begin{align*}
\frac{\partial^2}{\partial t^2}-\frac{\partial^2}{\partial x^2}+\frac{\partial}{\partial t}+\frac{\partial}{\partial x}
=\left(\frac{\partial}{\partial t}+\frac{\partial}{\partial x}\right)\left(\frac{\partial}{\partial t}-\frac{\partial}{\partial x}+1\right)
\end{align*}
将方程化为两个一阶的方程。
\end{example}

\begin{example}
求解 Goursat 问题
\begin{align*}
\begin{cases}
u_{tt}-u_{xx}=0, & t>|x|, \\
u|_{t=-x}=\varphi(x), & x\leqslant 0, \\
u|_{t=x}=\psi(x), & x\geqslant 0,
\end{cases}
\end{align*}
其中 $\varphi(0)=\psi(0)$。如果 $\varphi(x)$ 在 $(-a,0]$ ($a>0$) 上给定，$\psi(x)$ 在 $[0,b]$ ($b>0$) 上给定，试指出此定解条件的决定区域。
\end{example}

\begin{example}
求解 Darboux 问题
\begin{align*}
\begin{cases}
u_{tt}-u_{xx}=0, & 0<x<t, \\
u|_{x=0}=\varphi(t), & t\geqslant 0, \\
u|_{x=t}=\psi(t), & t\geqslant 0,
\end{cases}
\end{align*}
其中 $\varphi(0)=\psi(0)$。如果 $\varphi(t),\psi(t)$ 都在 $[0,a]$ ($a>0$) 上给定，试指出此定解条件的决定区域。
\end{example}

\begin{example}
求解问题
\begin{align*}
\begin{cases}
u_{tt}-u_{xx}=0, & 0<x<t, \\
u|_{x=t}=\varphi(t), & t\geqslant 0, \\
u|_{x=0}=\psi(t), & t\geqslant 0.
\end{cases}
\end{align*}
如果 $\varphi(t),\psi(t)$ 都在 $[0,a]$ ($a>0$) 上给定，试指出此定解条件的决定区域。
\end{example}

\begin{example}
求解问题
\begin{align*}
\begin{cases}
u_{tt}-u_{xx}=0, & 0<x<t, \\
u|_{x=t}=\varphi(t), & t\geqslant 0, \\
(-u_x+u)|_{x=0}=\psi(t), & t\geqslant 0.
\end{cases}
\end{align*}
如果 $\varphi(t),\psi(t)$ 都在 $[0,a]$ ($a>0$) 上给定，试指出此定解条件的决定区域。
\end{example}

\begin{example}
在区域 $[0,+\infty)\times[0,+\infty)$ 上考虑一阶方程
\begin{align*}
u_t+au_x=0 \quad\text{或}\quad u_t-au_x=0
\end{align*}
满足初边值条件
\begin{align*}
u|_{x=0}=g(t),\quad t\geqslant 0,\qquad u|_{t=0}=\varphi(x),\quad x\geqslant 0
\end{align*}
的初边值问题，这里 $a>0$ 是常数。试问对于哪一个方程上述初边值条件提法正确？为什么？对提法正确的初边值问题求解，并说明对 $g(t),\varphi(x)$ 提什么条件，所得的解是属于 $C^1$ 的解。
\end{example}

\begin{example}
求解并证明初值问题
\begin{align*}
\begin{cases}
u_{tt}-a^2u_{xx}+bu=f(x,t), & (x,t)\in\mathbb{R}_+^2, \\
u|_{t=0}=\varphi(x), & x\in\mathbb{R}, \\
u_t|_{t=0}=\psi(x), & x\in\mathbb{R}
\end{cases}
\end{align*}
解的惟一性，其中 $a\in\mathbb{R}_+,\ b\in\mathbb{R}$ 为常数。
提示：考虑二维初值问题
\begin{align*}
\begin{cases}
\tilde{u}_{tt}-a^2(\tilde{u}_{xx}+\tilde{u}_{yy})=\tilde{f}(x,y,t), & (x,y,t)\in\mathbb{R}_+^3, \\
\tilde{u}|_{t=0}=\tilde{\varphi}(x,y), & (x,y)\in\mathbb{R}^2, \\
\tilde{u}_t|_{t=0}=\tilde{\psi}(x,y), & (x,y)\in\mathbb{R}^2,
\end{cases}
\end{align*}
其中 $\tilde{u}=\tilde{u}(x,y,t)=u(x,t)v(y)$，$\tilde{f}=f(x,t)v(y)$，$\tilde{\varphi}=\varphi(x)v(y)$，$\tilde{\psi}=\psi(x)v(y)$，辅助函数 $v(y)$ 满足 $v''(y)=-\frac{b}{a^2}v(y)$。
\end{example}

\begin{example}
求解并证明初值问题
\begin{align*}
\begin{cases}
u_{tt}-a^2u_{xx}+bu_t+cu_x+du=f(x,t), & (x,t)\in\mathbb{R}_+^2, \\
u|_{t=0}=\varphi(x), & x\in\mathbb{R}, \\
u_t|_{t=0}=\psi(x), & x\in\mathbb{R}
\end{cases}
\end{align*}
解的惟一性。
\end{example}

\begin{example}
设 $u\in C^2(\mathbb{R}_+^2)$ 是 Cauchy 初值问题
\begin{align*}
\begin{cases}
u_{tt}-a^2u_{xx}=0, & (x,t)\in\mathbb{R}_+^2, \\
u(x,0)=\varphi(x), & x\in\mathbb{R}, \\
u_t(x,0)=\psi(x), & x\in\mathbb{R}
\end{cases}
\end{align*}
的解，其中 $a$ 为正常数，$\varphi\in C^2(\mathbb{R}),\ \psi\in C^1(\mathbb{R})$。证明：
\begin{enumerate}[(1)]
\item 对于任意 $t\geqslant 0$，成立能量不等式
\begin{align*}
\int_{-\infty}^{+\infty}[u_t^2(x,t)+a^2u_x^2(x,t)]\,\mathrm{d}x
\leqslant \int_{-\infty}^{+\infty}[\psi^2(x)+a^2|\varphi'(x)|^2]\,\mathrm{d}x;
\end{align*}
\item 对于任意 $t\geqslant 0$，成立能量恒等式
\begin{align*}
\int_{-\infty}^{+\infty}[u_t^2(x,t)+a^2u_x^2(x,t)]\,\mathrm{d}x
= \int_{-\infty}^{+\infty}[\psi^2(x)+a^2|\varphi'(x)|^2]\,\mathrm{d}x.
\end{align*}
\end{enumerate}
\end{example}

\begin{example}
设函数 $u\in C^2(\mathbb{R}_+^2)$ 满足第 37 题的 Cauchy 初值问题且 $\varphi(x),\psi(x)\in C_0^\infty(\mathbb{R})$。记其动能和势能分别为
\begin{align*}
k(t)=\frac{1}{2}\int_{-\infty}^{+\infty}u_t^2(x,t)\,\mathrm{d}x,\qquad
p(t)=\frac{a^2}{2}\int_{-\infty}^{+\infty}u_x^2(x,t)\,\mathrm{d}x.
\end{align*}
证明：
\begin{enumerate}[(1)]
\item $k(t)+p(t)$ 是与 $t$ 无关的常数；
\item 当 $t$ 足够大时，$k(t)=p(t)$。
\end{enumerate}
提示：利用 D'Alembert 公式和附注 4.2 的表达式，则
\begin{align*}
u_t(x,t)=aF'(x+at)-aG'(x-at),\qquad
u_x(x,t)=F'(x+at)+G'(x-at),
\end{align*}
其中
\begin{align*}
F'=\frac{1}{2a}(a\varphi'+\psi),\qquad G'=\frac{1}{2a}(a\varphi'-\psi).
\end{align*}
\end{example}

\begin{example}
设函数 $u\in C^2(\mathbb{R}_+^4)$ 满足 Cauchy 初值问题
\begin{align*}
\begin{cases}
u_{tt}-a^2\Delta u=0, & (x,t)\in\mathbb{R}_+^4, \\
u(x,0)=\varphi(x), & x\in\mathbb{R}^3, \\
u_t(x,0)=\psi(x), & x\in\mathbb{R}^3,
\end{cases}
\end{align*}
其中 $a$ 为正常数，$\varphi(x),\psi(x)\in C_0^\infty(\mathbb{R}^3)$。证明：存在常数 $C$，使得对于所有的 $(x,t)\in\mathbb{R}_+^4$，成立
\begin{align*}
|u(x,t)|\leqslant \frac{C}{t}.
\end{align*}
（提示：利用 Kirchhoff 公式）
\end{example}

\begin{example}
设 $a,h$ 为正常数，$A_1,A_2$ 为常数。用分离变量法求解下列混合问题：
\begin{enumerate}[(1)]
\item $\begin{cases}
u_{tt}-a^2u_{xx}=0, & 0<x<\pi,\ t>0, \\
u(x,0)=\sin x,\quad u_t(x,0)=x(\pi-x), & 0\leqslant x\leqslant\pi, \\
u(0,t)=0,\quad u(\pi,t)=0, & t\geqslant 0;
\end{cases}$
\item $\begin{cases}
u_{tt}-a^2u_{xx}=0, & 0<x<l,\ t>0, \\
u(x,0)=x(x-2l),\quad u_t(x,0)=0, & 0\leqslant x\leqslant l, \\
u(0,t)=0,\quad u_x(l,t)=0, & t\geqslant 0;
\end{cases}$
\item $\begin{cases}
u_{tt}-a^2u_{xx}=0, & 0<x<l,\ t>0, \\
u(x,0)=\cos\frac{\pi x}{l},\quad u_t(x,0)=0, & 0\leqslant x\leqslant l, \\
u_x(0,t)=0,\quad u_x(l,t)=0, & t\geqslant 0;
\end{cases}$
\item $\begin{cases}
u_{tt}-a^2u_{xx}=0, & 0<x<l,\ t>0, \\
u(x,0)=1,\quad u_t(x,0)=1, & 0\leqslant x\leqslant l, \\
u_x(0,t)=0,\quad u_x(l,t)=0, & t\geqslant 0;
\end{cases}$
\item $\begin{cases}
u_{tt}-u_{xx}=0, & 0<x<\pi,\ t>0, \\
u(x,0)=0,\quad u_t(x,0)=0, & 0\leqslant x\leqslant\pi, \\
u_x(0,t)=A_1t,\quad u_x(\pi,t)=A_2t, & t\geqslant 0;
\end{cases}$
\item $\begin{cases}
u_{tt}-a^2u_{xx}=0, & 0<x<l,\ t>0, \\
u(x,0)=0,\quad u_t(x,0)=\cos\frac{\pi}{l}x, & 0\leqslant x\leqslant l, \\
u_x|_{x=0}=0,\quad (u_x+hu)|_{x=l}=0, & t\geqslant 0.
\end{cases}$
\end{enumerate}
\end{example}
\begin{solution}
\begin{enumerate}[(1)]
\item 采用分离变量法。设非零解具有形式
\begin{align*}
u(x,t)=X(x)T(t),
\end{align*}
其中 $X\not\equiv 0$，$T\not\equiv 0$. 代入方程 $u_{tt}-a^2u_{xx}=0$，得
\begin{align*}
T''(t)X(x)-a^2X''(x)T(t)=0,\qquad (x,t)\in(0,\pi)\times(0,+\infty),
\end{align*}
即
\begin{align*}
\frac{T''(t)}{a^2T(t)}=\frac{X''(x)}{X(x)}.
\end{align*}
上式左端仅与 $t$ 有关，右端仅与 $x$ 有关，故必为常数，记为 $-\lambda$，于是
\begin{align*}
T''(t)+a^2\lambda T(t)=0,\qquad X''(x)+\lambda X(x)=0.
\end{align*}
将 $u=X(x)T(t)$ 代入边值条件 $u(0,t)=u(\pi,t)=0$，得
\begin{align*}
X(0)T(t)=X(\pi)T(t)=0,\qquad t>0.
\end{align*}
因 $T(t)\not\equiv 0$，故 $X(0)=X(\pi)=0$. 从而得到特征值问题
\begin{align*}
\begin{cases}
X''(x)+\lambda X(x)=0,\quad 0<x<\pi,\\
X(0)=X(\pi)=0.
\end{cases}
\end{align*}
求解该问题。若 $\lambda\leqslant 0$，则结合边值条件只能得到零解，与 $X\not\equiv 0$ 矛盾，故 $\lambda>0$. 令 $\mu=\sqrt{\lambda}$，则通解为
\begin{align*}
X(x)=C_1\cos \mu x+C_2\sin \mu x.
\end{align*}
由 $X(0)=0$ 得 $C_1=0$；由 $X(\pi)=0$ 得 $C_2\sin \mu\pi=0$. 因 $C_2\neq 0$，所以 $\sin \mu\pi=0$，即 $\mu=n$，$n=1,2,\cdots$. 故特征值和对应的特征函数为
\begin{align*}
\lambda_n=n^2,\qquad X_n(x)=\sin(nx),\quad n=1,2,\cdots.
\end{align*}
对于每个 $n$，方程 $T''+a^2\lambda_n T=0$ 的通解为
\begin{align*}
T_n(t)=A_n\cos(an t)+B_n\sin(an t),
\end{align*}
其中 $A_n,B_n$ 为任意常数。于是得到满足方程和边值条件的一系列特解
\begin{align*}
u_n(x,t)=\left(A_n\cos(an t)+B_n\sin(an t)\right)\sin(nx),\quad n=1,2,\cdots.
\end{align*}
由叠加原理，形式上构造
\begin{align*}
u(x,t)=\sum_{n=1}^{\infty}u_n(x,t)
=\sum_{n=1}^{\infty}\left(A_n\cos(an t)+B_n\sin(an t)\right)\sin(nx).
\end{align*}
为使 $u$ 满足初始条件，令
\begin{align}\label{eq:init1a}
u(x,0)=\varphi(x)=\sum_{n=1}^{\infty}A_n\sin(nx),\qquad
u_t(x,0)=\psi(x)=\sum_{n=1}^{\infty}anB_n\sin(nx),
\end{align}
其中 $\varphi(x)=\sin x$，$\psi(x)=x(\pi-x)$.
对于 $\varphi(x)=\sin x$，它本身就是 $n=1$ 的基函数，故
\begin{align*}
A_1=1,\qquad A_n=0\quad(n\geqslant 2).
\end{align*}
对于 $\psi(x)=x(\pi-x)$，其 Fourier 正弦系数为
\begin{align*}
anB_n=\frac{2}{\pi}\int_0^\pi x(\pi-x)\sin(nx)\,\mathrm{d}x,
\end{align*}
因此
\begin{align*}
B_n=\frac{2}{an\pi}\int_0^\pi x(\pi-x)\sin(nx)\,\mathrm{d}x.
\end{align*}
又因为
\begin{align*}
\int_0^\pi x(\pi-x)\sin(nx)\,\mathrm{d}x = \frac{2(1-(-1)^n)}{n^3},
\end{align*}
所以
\begin{align*}
B_n=\frac{4(1-(-1)^n)}{a\pi n^4}.
\end{align*}
当 $n$ 为偶数时 $B_n=0$；当 $n=2k-1$ 时，
\begin{align*}
B_{2k-1}=\frac{8}{a\pi(2k-1)^4}.
\end{align*}
综上可得
\begin{align*}
u(x,t)=\cos(at)\sin x+\sum_{k=1}^{\infty}\frac{8}{a\pi(2k-1)^4}\sin(a(2k-1)t)\sin((2k-1)x).
\end{align*}

\item 采用分离变量法。设非零解具有形式 $u(x,t)=X(x)T(t)$，代入方程 $u_{tt}-a^2u_{xx}=0$，得
\begin{align*}
\frac{T''(t)}{a^2T(t)}=\frac{X''(x)}{X(x)}=-\lambda,
\end{align*}
于是
\begin{align*}
T''(t)+a^2\lambda T(t)=0,\qquad X''(x)+\lambda X(x)=0.
\end{align*}
将 $u=X(x)T(t)$ 代入边值条件 $u(0,t)=0$，$u_x(l,t)=0$，得
\begin{align*}
X(0)=0,\qquad X'(l)=0.
\end{align*}
从而得到特征值问题
\begin{align*}
\begin{cases}
X''(x)+\lambda X(x)=0,\quad 0<x<l,\\
X(0)=0,\quad X'(l)=0.
\end{cases}
\end{align*}
求解该问题。若 $\lambda\leqslant 0$，则结合边值条件只能得到零解，与 $X\not\equiv 0$ 矛盾，故 $\lambda>0$. 令 $\mu=\sqrt{\lambda}$，则通解为
\begin{align*}
X(x)=C_1\cos \mu x+C_2\sin \mu x.
\end{align*}
由 $X(0)=0$ 得 $C_1=0$；由 $X'(l)=0$ 得 $C_2\mu\cos \mu l=0$. 因 $C_2\neq 0$，$\mu\neq 0$，所以 $\cos \mu l=0$，即
\begin{align*}
\mu_n=\frac{(2n-1)\pi}{2l},\qquad n=1,2,\cdots.
\end{align*}
故特征值和对应的特征函数为
\begin{align*}
\lambda_n=\mu_n^2=\left(\frac{(2n-1)\pi}{2l}\right)^2,\qquad X_n(x)=\sin(\mu_n x),\quad n=1,2,\cdots.
\end{align*}
对于每个 $n$，方程 $T''+a^2\lambda_n T=0$ 的通解为
\begin{align*}
T_n(t)=A_n\cos(a\mu_n t)+B_n\sin(a\mu_n t),
\end{align*}
其中 $A_n,B_n$ 为任意常数。由叠加原理，构造
\begin{align*}
u(x,t)=\sum_{n=1}^{\infty}\left(A_n\cos(a\mu_n t)+B_n\sin(a\mu_n t)\right)\sin(\mu_n x).
\end{align*}
由初速度 $u_t(x,0)=0$，得 $\sum a\mu_n B_n\sin(\mu_n x)=0$，故 $B_n=0$（$n=1,2,\cdots$）.
由初位移 $u(x,0)=x(x-2l)$，得
\begin{align*}
x(x-2l)=\sum_{n=1}^{\infty}A_n\sin(\mu_n x).
\end{align*}
因此
\begin{align*}
A_n=\frac{2}{l}\int_0^l x(x-2l)\sin(\mu_n x)\,\mathrm{d}x.
\end{align*}
又因为
\begin{align*}
\int_0^l (x^2-2lx)\sin(\mu_n x)\,\mathrm{d}x=-\frac{2}{\mu_n^3},
\end{align*}
所以
\begin{align*}
A_n=-\frac{4}{l\mu_n^3}=-\frac{32l^2}{(2n-1)^3\pi^3}.
\end{align*}
综上可得
\begin{align*}
u(x,t)=\sum_{n=1}^{\infty}\left[-\frac{32l^2}{(2n-1)^3\pi^3}\right]\cos(a\mu_n t)\sin(\mu_n x),
\end{align*}
其中 $\mu_n=\dfrac{(2n-1)\pi}{2l}$.

\item 采用分离变量法。设非零解具有形式 $u(x,t)=X(x)T(t)$，代入方程 $u_{tt}-a^2u_{xx}=0$，得
\begin{align*}
\frac{T''(t)}{a^2T(t)}=\frac{X''(x)}{X(x)}=-\lambda,
\end{align*}
于是
\begin{align*}
T''(t)+a^2\lambda T(t)=0,\qquad X''(x)+\lambda X(x)=0.
\end{align*}
将 $u=X(x)T(t)$ 代入边值条件 $u_x(0,t)=u_x(l,t)=0$，得
\begin{align*}
X'(0)=0,\qquad X'(l)=0.
\end{align*}
从而得到特征值问题
\begin{align*}
\begin{cases}
X''(x)+\lambda X(x)=0,\quad 0<x<l,\\
X'(0)=0,\quad X'(l)=0.
\end{cases}
\end{align*}
求解该问题。当 $\lambda=0$ 时，$X''=0$，通解 $X(x)=C_1+C_2x$，由 $X'(0)=X'(l)=0$ 得 $C_2=0$，故 $X_0(x)=1$（常数）.
当 $\lambda>0$ 时，令 $\mu=\sqrt{\lambda}$，通解为
\begin{align*}
X(x)=C_1\cos \mu x+C_2\sin \mu x.
\end{align*}
由 $X'(0)=0$ 得 $C_2=0$；由 $X'(l)=0$ 得 $-C_1\mu\sin \mu l=0$. 因 $C_1\neq 0$，$\mu\neq 0$，所以 $\sin \mu l=0$，即 $\mu_n=n\pi/l$，$n=1,2,\cdots$. 故特征值和对应的特征函数为
\begin{align*}
\lambda_0=0,\quad X_0(x)=1,\qquad
\lambda_n=\left(\frac{n\pi}{l}\right)^2,\quad X_n(x)=\cos\frac{n\pi}{l}x,\quad n=1,2,\cdots.
\end{align*}
对于 $\lambda_0=0$，方程 $T''=0$ 的通解为 $T_0(t)=A_0+B_0 t$.
对于 $n\geqslant 1$，方程 $T''+a^2\lambda_n T=0$ 的通解为
\begin{align*}
T_n(t)=A_n\cos\frac{na\pi}{l}t+B_n\sin\frac{na\pi}{l}t.
\end{align*}
由叠加原理，构造
\begin{align*}
u(x,t)=\frac{A_0+B_0t}{2}+\sum_{n=1}^{\infty}\left(A_n\cos\frac{na\pi}{l}t+B_n\sin\frac{na\pi}{l}t\right)\cos\frac{n\pi}{l}x.
\end{align*}
由初速度 $u_t(x,0)=0$，得
\begin{align*}
\frac{B_0}{2}+\sum_{n=1}^{\infty}\frac{na\pi}{l}B_n\cos\frac{n\pi}{l}x=0.
\end{align*}
故 $B_0=0$，$B_n=0$（$n\geqslant 1$）.
由初位移 $u(x,0)=\cos(\pi x/l)$，得
\begin{align*}
\frac{A_0}{2}+\sum_{n=1}^{\infty}A_n\cos\frac{n\pi}{l}x=\cos\frac{\pi x}{l}.
\end{align*}
该函数本身就是 $n=1$ 的余弦基函数，故 $A_0=0$，$A_1=1$，$A_n=0$（$n\geqslant 2$）.
综上可得
\begin{align*}
u(x,t)=\cos\frac{\pi x}{l}\cos\frac{a\pi t}{l}.
\end{align*}

\item 采用分离变量法。设非零解具有形式
\begin{align*}
u(x,t)=X(x)T(t),
\end{align*}
其中 $X\not\equiv 0$，$T\not\equiv 0$. 代入方程 $u_{tt}-a^2u_{xx}=0$，得
\begin{align*}
T''(t)X(x)-a^2X''(x)T(t)=0,\qquad (x,t)\in(0,l)\times(0,+\infty),
\end{align*}
即
\begin{align*}
\frac{T''(t)}{a^2T(t)}=\frac{X''(x)}{X(x)}.
\end{align*}
上式左端仅与 $t$ 有关，右端仅与 $x$ 有关，故必为常数，记为 $-\lambda$，于是
\begin{align*}
T''(t)+a^2\lambda T(t)=0,\qquad X''(x)+\lambda X(x)=0.
\end{align*}
将 $u=X(x)T(t)$ 代入边值条件 $u_x(0,t)=u_x(l,t)=0$，得
\begin{align*}
X'(0)T(t)=X'(l)T(t)=0,\qquad t>0.
\end{align*}
因 $T(t)\not\equiv 0$，故 $X'(0)=X'(l)=0$. 从而得到特征值问题
\begin{align*}
\begin{cases}
X''(x)+\lambda X(x)=0,\quad 0<x<l,\\
X'(0)=0,\quad X'(l)=0.
\end{cases}
\end{align*}
求解该问题。当 $\lambda=0$ 时，$X''=0$，通解 $X(x)=C_1+C_2x$，由 $X'(0)=X'(l)=0$ 得 $C_2=0$，故 $X_0(x)=1$（常数）.
当 $\lambda>0$ 时，令 $\mu=\sqrt{\lambda}$，通解为
\begin{align*}
X(x)=C_1\cos \mu x+C_2\sin \mu x.
\end{align*}
由 $X'(0)=0$ 得 $C_2=0$；由 $X'(l)=0$ 得 $-C_1\mu\sin \mu l=0$. 因 $C_1\neq 0$，$\mu\neq 0$，所以 $\sin \mu l=0$，即 $\mu_n=\dfrac{n\pi}{l}$，$n=1,2,\cdots$. 故特征值和对应的特征函数为
\begin{align*}
\lambda_0=0,\quad X_0(x)=1,\qquad
\lambda_n=\left(\frac{n\pi}{l}\right)^2,\quad X_n(x)=\cos\frac{n\pi}{l}x,\quad n=1,2,\cdots.
\end{align*}
对于 $\lambda_0=0$，方程 $T''=0$ 的通解为 $T_0(t)=A_0+B_0 t$.
对于 $n\geqslant 1$，方程 $T''+a^2\lambda_n T=0$ 的通解为
\begin{align*}
T_n(t)=A_n\cos\frac{na\pi}{l}t+B_n\sin\frac{na\pi}{l}t.
\end{align*}
于是得到满足方程和边值条件的一系列特解
\begin{align*}
u_0(x,t)=\frac{A_0+B_0t}{2},\qquad
u_n(x,t)=\left(A_n\cos\frac{na\pi}{l}t+B_n\sin\frac{na\pi}{l}t\right)\cos\frac{n\pi}{l}x,\quad n\geqslant 1.
\end{align*}
由叠加原理，构造
\begin{align*}
u(x,t)=\frac{A_0+B_0t}{2}+\sum_{n=1}^{\infty}\left(A_n\cos\frac{na\pi}{l}t+B_n\sin\frac{na\pi}{l}t\right)\cos\frac{n\pi}{l}x.
\end{align*}
为使 $u$ 满足初始条件，令
\begin{align*}
u(x,0)=1=\frac{A_0}{2}+\sum_{n=1}^{\infty}A_n\cos\frac{n\pi}{l}x,\qquad
u_t(x,0)=1=\frac{B_0}{2}+\sum_{n=1}^{\infty}\frac{na\pi}{l}B_n\cos\frac{n\pi}{l}x.
\end{align*}
常数 $1$ 的余弦展开只有常数项，故
\begin{align*}
\frac{A_0}{2}=1,\quad A_n=0\ (n\geqslant 1),\qquad
\frac{B_0}{2}=1,\quad B_n=0\ (n\geqslant 1).
\end{align*}
因此 $A_0=2$，$B_0=2$，$A_n=B_n=0$（$n\geqslant 1$）.
代入得
\begin{align*}
u(x,t)=\frac{2+2t}{2}=1+t.
\end{align*}

\item 边界条件为 $u_x(0,t)=A_1t$，$u_x(\pi,t)=A_2t$，非齐次，需先齐次化。寻找辅助函数 $\Phi(x)$ 满足
\begin{align*}
\Phi'(0)=A_1,\qquad \Phi'(\pi)=A_2.
\end{align*}
取 $\Phi(x)=\alpha x^2+\beta x$，则 $\Phi'(x)=2\alpha x+\beta$. 由 $\beta=A_1$，$2\alpha\pi+A_1=A_2$，得
\begin{align*}
\alpha=\frac{A_2-A_1}{2\pi},\qquad \beta=A_1.
\end{align*}
因此
\begin{align*}
\Phi(x)=\frac{A_2-A_1}{2\pi}x^2+A_1x.
\end{align*}
令 $u(x,t)=t\Phi(x)+v(x,t)$，则 $v$ 满足
\begin{align*}
v_x(0,t)=0,\qquad v_x(\pi,t)=0,
\end{align*}
且满足非齐次方程及初值：
\begin{align*}
v_{tt}-v_{xx}=t\frac{A_2-A_1}{\pi},\qquad
v(x,0)=0,\qquad v_t(x,0)=-\Phi(x).
\end{align*}
现在对 $v$ 用分离变量法。设非零解具有形式 $v(x,t)=X(x)T(t)$，代入齐次方程 $v_{tt}-v_{xx}=0$得
\begin{align*}
\frac{T''(t)}{T(t)}=\frac{X''(x)}{X(x)}=-\lambda,
\end{align*}
从而
\begin{align*}
T''(t)+\lambda T(t)=0,\qquad X''(x)+\lambda X(x)=0.
\end{align*}
边界条件 $v_x(0,t)=v_x(\pi,t)=0$ 给出 $X'(0)=X'(\pi)=0$. 当 $\lambda=0$ 时，$X_0(x)=1$；当 $\lambda>0$ 时，令 $\mu=\sqrt{\lambda}$，通解 $X(x)=C_1\cos \mu x+C_2\sin \mu x$，由 $X'(0)=0$ 得 $C_2=0$，由 $X'(\pi)=0$ 得 $-C_1\mu\sin \mu\pi=0$，故 $\mu=n$，$n=1,2,\cdots$. 因此特征值和特征函数为
\begin{align*}
\lambda_0=0,\quad X_0(x)=1,\qquad
\lambda_n=n^2,\quad X_n(x)=\cos(nx),\quad n=1,2,\cdots.
\end{align*}
将 $v$ 按此特征函数系展开：
\begin{align*}
v(x,t)=V_0(t)+\sum_{n=1}^{\infty}V_n(t)\cos(nx).
\end{align*}
代入非齐次方程 $v_{tt}-v_{xx}=t\dfrac{A_2-A_1}{\pi}$，得
\begin{align*}
V_0''(t)=t\frac{A_2-A_1}{\pi},\qquad
V_n''(t)+n^2V_n(t)=0\quad(n\geqslant 1).
\end{align*}
由初值条件 $v(x,0)=0$，$v_t(x,0)=-\Phi(x)$，得
\begin{align*}
V_0(0)=0,\qquad V_n(0)=0\quad(n\geqslant 1),
\end{align*}
\begin{align*}
V_0'(0)=-\frac{1}{\pi}\int_0^\pi\Phi(x)\,\mathrm{d}x,\qquad
V_n'(0)=-\frac{2}{\pi}\int_0^\pi\Phi(x)\cos(nx)\,\mathrm{d}x.
\end{align*}
计算积分：
\begin{align*}
\int_0^\pi\Phi(x)\,\mathrm{d}x=\frac{\pi^2(A_2+2A_1)}{6},
\quad
\int_0^\pi\Phi(x)\cos(nx)\,\mathrm{d}x
=\frac{A_2(-1)^n-A_1}{n^2}-\frac{A_2-A_1}{\pi n^3}.
\end{align*}
故
\begin{align*}
V_0'(0)=-\frac{\pi(A_2+2A_1)}{6},
\quad
V_n'(0)=-\frac{2}{\pi}\left[\frac{A_2(-1)^n-A_1}{n^2}-\frac{A_2-A_1}{\pi n^3}\right].
\end{align*}
解得
\begin{align*}
V_0(t)=\frac{A_2-A_1}{6\pi}t^3-\frac{\pi(A_2+2A_1)}{6}t.
\end{align*}
解 $V_n$（$n\geqslant 1$）：$V_n''+n^2V_n=0$，且 $V_n(0)=0$，故
\begin{align*}
V_n(t)=\frac{V_n'(0)}{n}\sin(nt).
\end{align*}
因此
\begin{align*}
V_n(t)=-\frac{2}{\pi n^3}\left(A_2(-1)^n-A_1\right)\sin(nt)+\frac{2(A_2-A_1)}{\pi^2n^4}\sin(nt).
\end{align*}
再由$u(x,t)=t\Phi(x)+v(x,t)$得
\begin{align*}
u(x,t)=t\left[\frac{A_2-A_1}{2\pi}x^2+A_1x\right]
+\frac{A_2-A_1}{6\pi}t^3-\frac{\pi(A_2+2A_1)}{6}t
+\sum_{n=1}^{\infty}\left[
-\frac{2(A_2(-1)^n-A_1)}{\pi n^3}
+\frac{2(A_2-A_1)}{\pi^2n^4}
\right]\sin(nt)\cos(nx).
\end{align*}

\item 采用分离变量法。设非零解具有形式 $u(x,t)=X(x)T(t)$，代入方程 $u_{tt}-a^2u_{xx}=0$，得
\begin{align*}
\frac{T''(t)}{a^2T(t)}=\frac{X''(x)}{X(x)}=-\lambda,
\end{align*}
于是
\begin{align*}
T''(t)+a^2\lambda T(t)=0,\qquad X''(x)+\lambda X(x)=0.
\end{align*}
将 $u=X(x)T(t)$ 代入边值条件 $u_x(0,t)=0$，$(u_x+hu)(l,t)=0$，得
\begin{align*}
X'(0)=0,\qquad X'(l)+hX(l)=0.
\end{align*}
从而得到特征值问题
\begin{align*}
\begin{cases}
X''(x)+\lambda X(x)=0,\quad 0<x<l,\\
X'(0)=0,\quad X'(l)+hX(l)=0.
\end{cases}
\end{align*}
求解该问题。若 $\lambda\leqslant 0$，则结合边值条件只能得到零解，与 $X\not\equiv 0$ 矛盾，故 $\lambda>0$. 令 $\mu=\sqrt{\lambda}$，则通解为
\begin{align*}
X(x)=C_1\cos \mu x+C_2\sin \mu x.
\end{align*}
由 $X'(0)=0$ 得 $C_2=0$，故 $X(x)=C_1\cos \mu x$.
代入右端边界条件：
\begin{align*}
X'(l)+hX(l)=-C_1\mu\sin \mu l+hC_1\cos \mu l=0.
\end{align*}
因 $C_1\neq 0$，得
\begin{align*}
\mu\tan(\mu l)=h.
\end{align*}
设该超越方程的所有正根为
\begin{align*}
0<\mu_1<\mu_2<\cdots<\mu_n<\cdots,
\end{align*}
则特征值和对应的特征函数为
\begin{align*}
\lambda_n=\mu_n^2,\qquad X_n(x)=\cos(\mu_n x),\quad n=1,2,\cdots.
\end{align*}
对于每个 $n$，方程 $T''+a^2\lambda_n T=0$ 的通解为
\begin{align*}
T_n(t)=A_n\cos(a\mu_n t)+B_n\sin(a\mu_n t).
\end{align*}
由叠加原理，构造
\begin{align*}
u(x,t)=\sum_{n=1}^{\infty}\left(A_n\cos(a\mu_n t)+B_n\sin(a\mu_n t)\right)\cos(\mu_n x).
\end{align*}
由初位移 $u(x,0)=0$，得 $A_n=0$（$n=1,2,\cdots$）.
由初速度 $u_t(x,0)=\cos(\pi x/l)$，得
\begin{align*}
\cos\frac{\pi x}{l}=\sum_{n=1}^{\infty}a\mu_n B_n\cos(\mu_n x).
\end{align*}
因此
\begin{align*}
a\mu_n B_n=C_n,
\end{align*}
其中
\begin{align*}
C_n=\frac{\int_0^l\cos(\pi x/l)\cos(\mu_n x)\,\mathrm{d}x}{\int_0^l\cos^2(\mu_n x)\,\mathrm{d}x}.
\end{align*}
计算分母：
\begin{align*}
\int_0^l\cos^2(\mu_n x)\,\mathrm{d}x=\frac{l}{2}+\frac{\sin(2\mu_n l)}{4\mu_n}.
\end{align*}
计算分子：
\begin{align*}
\int_0^l\cos(\pi x/l)\cos(\mu_n x)\,\mathrm{d}x
=\frac{\mu_n\sin(\mu_n l)}{\left(\frac{\pi}{l}\right)^2-\mu_n^2}.
\end{align*}
综上可得
\begin{align*}
u(x,t)=\sum_{n=1}^{\infty}\frac{C_n}{a\mu_n}\sin(a\mu_n t)\cos(\mu_n x),
\end{align*}
其中 $\mu_n$ 是方程 $\mu\tan(\mu l)=h$ 的正根，且
\begin{align*}
C_n=\frac{\dfrac{\mu_n\sin(\mu_n l)}{\left(\frac{\pi}{l}\right)^2-\mu_n^2}}{\dfrac{l}{2}+\dfrac{\sin(2\mu_n l)}{4\mu_n}}.
\end{align*}
\end{enumerate}
\end{solution}

\begin{example}
设 $a$ 为正常数，$A,\omega,A_1$ 为常数。用分离变量法求解下列混合问题：
\begin{enumerate}[(1)]
\item $\begin{cases}
u_{tt}-a^2u_{xx}=\sin\frac{\pi}{l}x, & 0<x<l,\ t>0, \\
u(x,0)=0,\quad u_t(x,0)=0, & 0\leqslant x\leqslant l, \\
u(0,t)=0,\quad u(l,t)=0, & t\geqslant 0;
\end{cases}$
\item $\begin{cases}
u_{tt}-a^2u_{xx}+2u_t=x(l-x), & 0<x<l,\ t>0, \\
u(x,0)=0,\quad u_t(x,0)=0, & 0\leqslant x\leqslant l, \\
u(0,t)=0,\quad u(l,t)=0, & t\geqslant 0;
\end{cases}$
\item $\begin{cases}
u_{tt}-a^2u_{xx}=0, & 0<x<l,\ t>0, \\
u(x,0)=\frac{Ax^2}{l^2},\quad u_t(x,0)=0, & 0\leqslant x\leqslant l, \\
u_x(0,t)=0,\quad u(l,t)=A, & t\geqslant 0;
\end{cases}$
\item $\begin{cases}
u_{tt}-a^2u_{xx}=0, & 0<x<l,\ t>0, \\
u(x,0)=1,\quad u_t(x,0)=0, & 0\leqslant x\leqslant l, \\
u_x(0,t)=A\sin\omega t,\quad u(l,t)=1, & t\geqslant 0,
\end{cases}$
分别讨论产生共振和不产生共振两种情形；
\item $\begin{cases}
u_{tt}-u_{xx}=0, & 0<x<l,\ t>0, \\
u(x,0)=0,\quad u_t(x,0)=0, & 0\leqslant x\leqslant l, \\
u_x|_{x=0}=A_1t,\quad (u_x+u)|_{x=l}=0, & t\geqslant 0.
\end{cases}$
\end{enumerate}
\end{example}

\begin{example}
求解混合问题
\begin{align*}
\begin{cases}
u_{tt}-u_{xx}-\frac{2}{x}u_x=0, & 0<x<1,\ t>0, \\
u(x,0)=0,\quad u_t(x,0)=\cos\frac{\pi}{2}x, & 0\leqslant x\leqslant 1, \\
u(0,t)=u(1,t)=0, & t\geqslant 0.
\end{cases}
\end{align*}
（提示：函数 $v(x,t)=xu(x,t)$ 满足 $v_{tt}-v_{xx}=0$）
\end{example}

\begin{example}
用分离变量法求解混合问题
\begin{align*}
u_{tt}-(u_{xx}+u_{yy})=0,\quad 0<x<1,\ 0<y<1,\ t>0,
\end{align*}
\begin{align*}
u|_{t=0}=x(1-x),\quad 0\leqslant x\leqslant 1,\ 0\leqslant y\leqslant 1,
\end{align*}
\begin{align*}
u_t|_{t=0}=y(1-y),\quad 0\leqslant x\leqslant 1,\ 0\leqslant y\leqslant 1,
\end{align*}
\begin{align*}
u|_{x=0}=u|_{x=1}=0,\quad 0\leqslant y\leqslant 1,\ t\geqslant 0,
\end{align*}
\begin{align*}
u_y|_{y=0}=u_y|_{y=1}=0,\quad 0\leqslant x\leqslant 1,\ t\geqslant 0.
\end{align*}
\end{example}

\begin{example}
设 $u(x,t)$ 满足混合问题
\begin{align*}
\begin{cases}
u_{tt}-a^2u_{xx}=f(x,t), & 0<x<l,\ t>0, \\
u(x,0)=\varphi(x),\quad u_t(x,0)=\psi(x), & 0\leqslant x\leqslant l, \\
(-u_x+\alpha u)|_{x=0}=g_1(t), & t\geqslant 0, \\
(u_x+\beta u)|_{x=l}=g_2(t), & t\geqslant 0,
\end{cases}
\end{align*}
其中 $a$ 为正常数。试引进辅助函数，将边值条件齐次化。设
\begin{enumerate}[(1)]
\item $\alpha>0,\ \beta>0$；
\item $\alpha=\beta=0$。
\end{enumerate}
\end{example}

在以下问题中记 $Q_T=(0,l)\times(0,T)$。

\begin{example}
设 $u(x,t)\in C^1(\overline{Q}_T)\cap C^2(\overline{Q}_T)$ 满足混合问题
\begin{align*}
\begin{cases}
u_{tt}-a^2u_{xx}=f(x,t), & (x,t)\in Q_T, \\
u(x,0)=\varphi(x),\quad u_t(x,0)=\psi(x), & 0\leqslant x\leqslant l, \\
u_x(0,t)=0,\quad u_x(l,t)=0, & 0\leqslant t\leqslant T,
\end{cases}
\end{align*}
其中 $a$ 为正常数。试推导能量不等式。
\end{example}

\begin{example}
设 $u(x,t)\in C^1(\overline{Q}_T)\cap C^2(\overline{Q}_T)$ 满足混合问题
\begin{align*}
\begin{cases}
u_{tt}-a^2u_{xx}=f(x,t), & (x,t)\in Q_T, \\
u(x,0)=\varphi(x),\quad u_t(x,0)=\psi(x), & 0\leqslant x\leqslant l, \\
(-u_x+\alpha u)|_{x=0}=0,\quad (u_x+\beta u)|_{x=l}=0, & 0\leqslant t\leqslant T,
\end{cases}
\end{align*}
其中 $a,\alpha,\beta>0$ 为常数。试推导能量不等式。
\end{example}

\begin{example}
设 $u(x,t)\in C^1(\overline{Q}_T)\cap C^2(\overline{Q}_T)$ 满足混合问题
\begin{align*}
\begin{cases}
u_{tt}-a^2u_{xx}=f(x,t), & (x,t)\in Q_T, \\
u(x,0)=\varphi(x),\quad u_t(x,0)=\psi(x), & 0\leqslant x\leqslant l, \\
u_x|_{x=0}=0,\quad (u_x+u)|_{x=l}=0, & 0\leqslant t\leqslant T,
\end{cases}
\end{align*}
其中 $a$ 为正常数。试推导能量不等式。
\end{example}

\begin{example}
考虑定解问题
\begin{align*}
\begin{cases}
u_{tt}-u_{xx}=f(x,t), & (x,t)\in Q_T, \\
u(x,0)=\varphi(x),\quad u_t(x,0)=\psi(x), & 0\leqslant x\leqslant l, \\
u(0,t)=u(l,t)=0, & 0\leqslant t\leqslant T.
\end{cases}
\end{align*}
试问：对 $\varphi,\psi,f$ 提什么条件才能保证由分离变量法所得的解是古典解？试证明之。
\end{example}

\begin{example}
试引入定解问题
\begin{align*}
\begin{cases}
u_{tt}-u_{xx}=f(x,t), & (x,t)\in Q_T, \\
u(x,0)=\varphi(x),\quad u_t(x,0)=\psi(x), & 0\leqslant x\leqslant l, \\
u(0,t)=u(l,t)=0, & 0\leqslant t\leqslant T
\end{cases}
\end{align*}
的广义解的定义。试问：对 $\varphi,\psi,f$ 提什么条件才能保证上述定解问题的广义解的存在性唯一性？试证明之。
\end{example}

\begin{example}
设 $u(x,t)\in C(\overline{Q}_T)$ 是混合问题 \eqref{eq:wave-gen-mixed} 的广义解。证明：
\begin{enumerate}[(1)]
\item $u(x,0)=\varphi(x),\quad x\in[0,l]$；
\item 如果 $u(x,t)\in C^1(\overline{Q}_T)$，则 $u_t(x,0)=\psi(x),\quad x\in[0,l]$。
\end{enumerate}
\end{example}




\end{document}